% 02_problema.tex — Slide 2: Definicion del problema
\begin{frame}{Problema: 3-Dimensional Matching (Karp \#17)}

\begin{columns}[T]
\begin{column}{0.52\textwidth}
\textbf{Datos:}
\begin{itemize}
  \item Tres conjuntos $X, Y, Z$ con $|X|{=}|Y|{=}|Z|{=}n$
  \item Conjunto de tripletas $T \subseteq X \times Y \times Z$, $|T|{=}m$
\end{itemize}

\medskip
\textbf{Objetivo (max-3DM):}
Encontrar el mayor subconjunto $M \subseteq T$ tal que cada elemento aparece a lo sumo una vez en $M$.

\medskip
\textbf{Ejemplo} ($n{=}3$, $m{=}5$):
\begin{itemize}
  \item $M^* = \{(x_1,y_1,z_1),\,(x_2,y_2,z_2),\,(x_3,y_3,z_3)\}$
  \item $|M^*| = 3 = n$ \alert{(matching perfecto)}
\end{itemize}
\end{column}

\begin{column}{0.44\textwidth}
\begin{tikzpicture}[scale=0.82, every node/.style={font=\small}]
  % Nodos X (izquierda)
  \foreach \i/\label in {1/$x_1$, 2/$x_2$, 3/$x_3$} {
    \node[draw, circle, fill=blue!15, minimum size=20pt] (x\i) at (0, -\i*1.1+1.6) {\label};
  }
  % Nodos Y (centro)
  \foreach \i/\label in {1/$y_1$, 2/$y_2$, 3/$y_3$} {
    \node[draw, circle, fill=green!15, minimum size=20pt] (y\i) at (2, -\i*1.1+1.6) {\label};
  }
  % Nodos Z (derecha)
  \foreach \i/\label in {1/$z_1$, 2/$z_2$, 3/$z_3$} {
    \node[draw, circle, fill=red!15, minimum size=20pt] (z\i) at (4, -\i*1.1+1.6) {\label};
  }
  % Aristas del matching optimo (gruesas, azules)
  \draw[thick, mainblue] (x1) -- (y1) -- (z1);
  \draw[thick, mainblue] (x2) -- (y2) -- (z2);
  \draw[thick, mainblue] (x3) -- (y3) -- (z3);
  % Aristas no usadas (finas, grises)
  \draw[thin, gray!50] (x1) -- (y2);
  \draw[thin, gray!50] (x1) to[bend left=15] (z3);
  \draw[thin, gray!50] (x3) -- (y1);
  % Etiquetas columnas
  \node[font=\bfseries\small, blue!70] at (0, 1.3) {$X$};
  \node[font=\bfseries\small, green!60!black] at (2, 1.3) {$Y$};
  \node[font=\bfseries\small, red!70] at (4, 1.3) {$Z$};
\end{tikzpicture}
\end{column}
\end{columns}

\end{frame}
